\SourceLeaf{401}%
oder
\[
\alpha=1,\quad\beta=0
\]
\par
Der Nenner in dem Werthe von \(\xi,\eta\) ist der doppelte Inhalt des Dreiecks.

\par\medskip\Needspace{5\baselineskip}
\pdfbookmark[2]{II.}{carnot-zusatz-ii}
\centerline{\textbf{II.}}\nobreak\smallskip
Dass die Perpendikel in einem Dreiecke, aus den Spitzen auf die gegenüberstehenden Seiten sich in einem Punkte schneiden, kann man sehr einfach so zeigen.

\begingroup\rightskip=0pt plus 1em
Das gegebene Dreieck sei \(BDF\), und die erwähnten Perpendikel \(\overline{BI},\overline{DG},\overline{FH}\).\par
\endgroup

Man ziehe durch jedem Scheitelpunkt des Dreiecks Parallelen mit der gegenüberstehenden Seite, die sich in den Punkten \(A,C,E\), schneiden, es steht folglich \(\overline{FH}\) auch auf \(\overline{AE}\), \(\overline{GD}\) auf \(\overline{CE}\), \(\overline{BI}\) auf \(\overline{AC}\) senkrecht, und zwar ist
\[
\begin{aligned}\overline{AB}&=\overline{BC}\\\overline{ED}&=\overline{DC}\\\overline{AF}&=\overline{FE}\end{aligned}
\]
Beschreibt man nun um das Dreieck \(ACE\) einen Kreis, so liegt sein Mittelpunkt sowohl in \(\overline{BI}\), als in \(\overline{DG}\), als in \(\overline{FH}\), diese drei Linien müssen sich also in einem Punkte schneiden.

\textsc{Puissant} gibt in seinem Recueil des propositions de Géométrie einen zierlichen analytischen Beweis, und fügt einen geometrischen bei, der nicht dasselbe Verdienst hat.

\par\medskip\Needspace{5\baselineskip}
\pdfbookmark[2]{Erste handschriftliche Bemerkung.}{carnot-handschriftlich-i}
\begin{center}[\textit{Erste handschriftliche Bemerkung.}]\end{center}
\hspace*{\parindent}Sind \(\alpha,\beta,\gamma\) die complexen Zahlen, die sich in natürlicher Ordnung auf die drei Winkelpunkte eines Dreiecks beziehen; \(A,B,C\) die drei Winkel, \(u\) die complexe Zahl für den Mittelpunkt des [umschriebenen] Kreises, so hat man
\[
\begin{aligned}
2u&=\alpha+\beta+(\beta-\alpha)\operatorname{cotg}C.i=\alpha(1-i\operatorname{cotg}C)+\beta(1+i\operatorname{cotg}C)\\
&\hphantom{=\alpha+\beta+(\beta-\alpha)\operatorname{cotg}C.i}=\beta(1-i\operatorname{cotg}A)+\gamma(1+i\operatorname{cotg}A)\\
&\hphantom{=\alpha+\beta+(\beta-\alpha)\operatorname{cotg}C.i}=\gamma(1-i\operatorname{cotg}B)+\alpha(1+i\operatorname{cotg}B)
\end{aligned}
\]
